% GENERATED FILE. Edit canonical lessons, then run npm run generate:books.
% canonical-source-hash: fnv1a64:4c027ba5734a215c
% canonical-lessons: SA-C94-bahih, SA-C94-upari, SA-C94-adhah, SA-C94-agre, SA-C94-prsthatah

\chapter{Outside, Above, Below, In front, Behind}
\label{ch:sa-outside-above-below-in-front-behind}
\hlchaptermodality{94}

\begin{chapteropening}
\textbf{By the end of this chapter:} \emph{I can say outside, above, below, in front and behind.}

Complete the last lesson of chapter 94: I can say outside, above, below, in front and behind.
\end{chapteropening}

\section[\texorpdfstring{\sk{बहिः} (bahiḥ)}{बहिः (bahiḥ)}]{\sk{बहिः} (bahiḥ) — outside}

\label{lesson:SA-C94-bahih}



\begin{quote}
Before the new one: say the Sanskrit for right; south, then the Sanskrit for inside.
\end{quote}

\subsection*{You'll want to know: \sk{बहिः}}
\textbf{\sk{बहिः}} — \emph{bahiḥ} — \textquotedblleft{}outside\textquotedblright{}.

\begin{grammarlens}[title={What you've built}]
One more word for this chapter.
\end{grammarlens}

\subsection*{Your turn}
\begin{itemize}
\raggedright
  \item \emph{Say it:} \emph{bahiḥ}
  \item \emph{Say it:} \emph{bahiḥ}, once more
  \item \emph{Say it:} say \emph{bahiḥ}
  \item \emph{Recall:} say the Sanskrit for right; south, then the Sanskrit for inside, then say \emph{bahiḥ} again
\end{itemize}

\begin{tcolorbox}[breakable,colback=teal!4,colframe=teal!35!black,title={Before you move on}]
What does \sk{बहिः} mean? (\textquotedblleft{}Outside\textquotedblright{}.) Say it once more.
\end{tcolorbox}
\section[\texorpdfstring{\sk{उपरि} (upari)}{उपरि (upari)}]{\sk{उपरि} (upari) — above}

\label{lesson:SA-C94-upari}



\begin{quote}
Before the new one: say the Sanskrit for inside, then the Sanskrit for outside.
\end{quote}

\subsection*{You'll want to know: \sk{उपरि}}
\textbf{\sk{उपरि}} — \emph{upari} — \textquotedblleft{}above\textquotedblright{}.

\begin{grammarlens}[title={What you've built}]
One more word for this chapter.
\end{grammarlens}

\subsection*{Your turn}
\begin{itemize}
\raggedright
  \item \emph{Say it:} \emph{upari}
  \item \emph{Say it:} \emph{upari}, once more
  \item \emph{Say it:} say \emph{upari}
  \item \emph{Recall:} say the Sanskrit for inside, then the Sanskrit for outside, then say \emph{upari} again
\end{itemize}

\begin{tcolorbox}[breakable,colback=teal!4,colframe=teal!35!black,title={Before you move on}]
What does \sk{उपरि} mean? (\textquotedblleft{}Above\textquotedblright{}.) Say it once more.
\end{tcolorbox}
\section[\texorpdfstring{\sk{अधः} (adhaḥ)}{अधः (adhaḥ)}]{\sk{अधः} (adhaḥ) — below}

\label{lesson:SA-C94-adhah}



\begin{quote}
Before the new one: say the Sanskrit for outside, then the Sanskrit for above.
\end{quote}

\subsection*{You'll want to know: \sk{अधः}}
\textbf{\sk{अधः}} — \emph{adhaḥ} — \textquotedblleft{}below\textquotedblright{}.

\begin{grammarlens}[title={What you've built}]
One more word for this chapter.
\end{grammarlens}

\subsection*{Your turn}
\begin{itemize}
\raggedright
  \item \emph{Say it:} \emph{adhaḥ}
  \item \emph{Say it:} \emph{adhaḥ}, once more
  \item \emph{Say it:} say \emph{adhaḥ}
  \item \emph{Recall:} say the Sanskrit for outside, then the Sanskrit for above, then say \emph{adhaḥ} again
\end{itemize}

\begin{tcolorbox}[breakable,colback=teal!4,colframe=teal!35!black,title={Before you move on}]
What does \sk{अधः} mean? (\textquotedblleft{}Below\textquotedblright{}.) Say it once more.
\end{tcolorbox}
\section[\texorpdfstring{\sk{अग्रे} (agre)}{अग्रे (agre)}]{\sk{अग्रे} (agre) — in front}

\label{lesson:SA-C94-agre}



\begin{quote}
Before the new one: say the Sanskrit for above, then the Sanskrit for below.
\end{quote}

\subsection*{You'll want to know: \sk{अग्रे}}
\textbf{\sk{अग्रे}} — \emph{agre} — \textquotedblleft{}in front\textquotedblright{}.

\begin{grammarlens}[title={What you've built}]
One more word for this chapter.
\end{grammarlens}

\subsection*{Your turn}
\begin{itemize}
\raggedright
  \item \emph{Say it:} \emph{agre}
  \item \emph{Say it:} \emph{agre}, once more
  \item \emph{Say it:} say \emph{agre}
  \item \emph{Recall:} say the Sanskrit for above, then the Sanskrit for below, then say \emph{agre} again
\end{itemize}

\begin{tcolorbox}[breakable,colback=teal!4,colframe=teal!35!black,title={Before you move on}]
What does \sk{अग्रे} mean? (\textquotedblleft{}In front\textquotedblright{}.) Say it once more.
\end{tcolorbox}
\section[\texorpdfstring{\sk{पृष्ठतः} (pṛṣṭhataḥ)}{पृष्ठतः (pṛṣṭhataḥ)}]{\sk{पृष्ठतः} (pṛṣṭhataḥ) — behind}

\label{lesson:SA-C94-prsthatah}



\begin{quote}
Before the new one: say the Sanskrit for below, then the Sanskrit for in front.
\end{quote}

\subsection*{You'll want to know: \sk{पृष्ठतः}}
\textbf{\sk{पृष्ठतः}} — \emph{pṛṣṭhataḥ} — \textquotedblleft{}behind\textquotedblright{}.

\begin{grammarlens}[title={What you've built}]
One more word for this chapter.
\end{grammarlens}

\subsection*{Your turn}
\begin{itemize}
\raggedright
  \item \emph{Say it:} \emph{pṛṣṭhataḥ}
  \item \emph{Say it:} \emph{pṛṣṭhataḥ}, once more
  \item \emph{Say it:} say \emph{pṛṣṭhataḥ}
  \item \emph{Recall:} say the Sanskrit for below, then the Sanskrit for in front, then say \emph{pṛṣṭhataḥ} again
\end{itemize}

\begin{tcolorbox}[breakable,colback=teal!4,colframe=teal!35!black,title={Before you move on}]
What does \sk{पृष्ठतः} mean? (\textquotedblleft{}Behind\textquotedblright{}.) Say it once more.
\end{tcolorbox}
